% !TeX program = xelatex
% UTF-8。使用 XeLaTeX 编译两次即可。参考文献内置，无需 BibTeX 或外部图片。
\documentclass[12pt,a4paper]{article}
\usepackage[UTF8,scheme=plain,heading=false,fontset=none]{ctex}
\IfFontExistsTF{SimSun}
  {\setCJKmainfont[BoldFont=SimHei]{SimSun}}
  {\IfFontExistsTF{Noto Serif CJK SC}
    {\setCJKmainfont{Noto Serif CJK SC}}
    {\setCJKmainfont[BoldFont=FandolSong-Bold.otf]{FandolSong-Regular.otf}}}
\usepackage[margin=1in]{geometry}
\usepackage{array,tabularx,booktabs,titlesec,setspace}
\usepackage[table]{xcolor}
\usepackage[numbers]{natbib}
\usepackage[unicode,hidelinks]{hyperref}
\hypersetup{pdftitle={分子 DFT 计算中四种 ORCA 工作流的对照研究},
  pdfauthor={CHEN GENG},pdflang={zh-CN}}
\singlespacing
\setlength{\parindent}{0pt}
\setlength{\parskip}{0.36em}
\setlength{\tabcolsep}{5pt}
\setlength{\emergencystretch}{1.2em}
\renewcommand{\arraystretch}{1.18}
\newcolumntype{Y}{>{\raggedright\arraybackslash}X}
\newcolumntype{P}[1]{>{\raggedright\arraybackslash}p{#1}}
\titleformat{\section}{\normalsize\bfseries}{\thesection.}{0.5em}{}
\titlespacing*{\section}{0pt}{0.80em}{0.32em}
\newcommand{\subhead}[1]{\par\addvspace{0.4em}\textbf{#1}\par\nobreak}
\newcommand{\kJmol}{\ensuremath{\mathrm{kJ\,mol^{-1}}}}
\newcommand{\doi}[1]{\href{https://doi.org/#1}{doi:\,#1}}
\newcommand{\blankfield}{\rule{3.5cm}{0.35pt}}
\newcommand{\StudentName}{CHEN GENG}
\newcommand{\StudentID}{\blankfield}
\newcommand{\Supervisor}{\blankfield}
\newcommand{\SubmissionDate}{\blankfield} % DD/MM/YYYY
\raggedbottom
\begin{document}

\textbf{BG6023 硕士研究 I}\hfill CA1 研究提案\par
\textbf{项目名称\quad 分子 DFT 计算中四种 ORCA 工作流的对照研究}\par
\begin{tabularx}{\linewidth}{|P{0.456\linewidth}|Y|}
\hline
\rowcolor[HTML]{F2F2F2}\textbf{项目} & \textbf{信息} \\ \hline
学生姓名 & \StudentName \\ \hline
学号 & \StudentID \\ \hline
导师 & \Supervisor \\ \hline
提交日期 & \SubmissionDate \\ \hline
\end{tabularx}

\section{引言与研究背景}
用密度泛函理论（DFT）回答分子问题，需要把题目转成输入、完成计算，并从输出中读出正确的物理量。Jacobs 与 Pollice 检验了大语言模型能否生成 ORCA 可执行的输入文件\cite{jacobs}。Q-Planner 让模型规划任务，由确定性程序执行\cite{qplanner}。Quntur 的研究则在相同底层模型下，将专用智能体与加载 ORCA 技能的通用智能体进行了比较\cite{quntur}。对已经拥有成熟计算流程的实验室，选用哪种工作流取决于题目。固定脚本能完成哪些任务？智能体承担规划后，又能节省多少人工操作时间？

作业结束后还须决定是否继续计算。同一对气相和溶剂中的单点能可以回答电子能差问题，推断水合自由能则还需要热力学条件。对柔性分子，CENSO 先以低成本方法筛选构象，再提高理论水平精算\cite{censo}。如果只能精算四个候选，前一次得到的 DFT 能量是否有助于选中下一个候选，可以通过保留完整参考能量表来检验。

\section{研究问题的依据}
我在前期 Auto-ORCA 测试中记录到，八次 ORCA 启动有七次正常结束，但七道任务均缺少至少一项必交结果。例如，丁烷扫描完成了六个点，却遗漏了指定端点。前期乙醇计算在两套基组下均得到负的水相减气相电子能差，六个正己烷几何的电子能跨度达 $12.61$~\kJmol{}。这些先导结果说明，有必要分别检验计算是否完成、输出能否支持结论，以及有限预算下会错过多低的能量。

\section{研究目的、目标与问题}
\textbf{总体目的}\quad 在相同分子任务上比较四种 ORCA 工作流的科学正确性与成本，并检验专用 Auto-ORCA 在计算结束后作出的两类决策。

\textbf{目标一\quad 比较四种工作流。}在六类共 24 道题中，各自动化流程有多少题能复现经独立复核的人工结果？各流程需要多少人工时间、计算时间和返工？

\textbf{目标二\quad 检验输出复用。}对 12 种分子各提出一对后续问题，显式核对目标物理量与计算条件，能否提高两问均答对的比例，并减少缺乏依据的数值回答？

\textbf{目标三\quad 评估构象选择。}对 10 种柔性分子，各给四次查询机会。逐次提供 DFT 能量是否提高选中最低能候选的概率，并缩小错过最低点的能量差？

\clearpage
\section{研究方法}
\subhead{研究一\quad 四种工作流共用的 24 道独立测试题}
每类选用四种不同分子，每种分子只在主评测中出现一次。下表列出拟冻结的题目。坐标、原子编号、构型 A/B、电子态和判分规则均在正式测试前固定。开发中用过的同题同坐标仅用于预试。四组接收相同的科学条件，两个智能体另接收逐字相同的自然语言题面。

\begin{tabularx}{\linewidth}{@{}P{0.19\linewidth}P{0.33\linewidth}Y@{}}
\toprule
\textbf{题型} & \textbf{每类四题，各用一种分子} & \textbf{必交结果与参考容差} \\
\midrule
溶剂电子能差 & 乙腈、二甲醚、乙酸甲酯、吡啶 & 同一几何计算水相减气相电子能。方向、单位正确，误差不超过 $0.1$~\kJmol{}。\\
双构型电子能差 & 1,2-二氟乙烷、1,2-二氯乙烷、1-氟丙烷、乙二醇 & 对给定的两个几何做气相单点。报告 $E_B-E_A$ 和较低者，容差 $0.1$~\kJmol{}。\\
优化几何与键长 & 乙烯、乙炔、氟甲烷、甲酰胺 & 核对优化收敛。分别测量指定的 C=C、C$\equiv$C、C--F、C--N 键，容差 $0.005$~\AA。\\
羰基振动归属 & 乙醛、环丁酮、甲酸甲酯、乙酰胺 & 优化后在同一几何计算频率，检查极小点，以振动位移归属 C=O 伸缩。容差 $2$~cm$^{-1}$。\\
垂直单重激发 & 1,3-丁二烯、丙烯醛、呋喃、吡嗪 & 给出六个激发能和长度表象振子强度，逐根能量容差 $0.002$~eV，识别最亮跃迁。\\
固定几何偶极矩 & 甲醇、乙胺、氯甲烷、二硫化碳 & 提取总偶极矩模长，容差 $0.01$~D。二硫化碳采用线性对称坐标以检验零值。\\
\bottomrule
\end{tabularx}

溶剂题采用 B3LYP/def2-SVP 与 CPCM 隐式水，其余五类采用气相 PBE0/def2-SVP。双构型题固定全部坐标，以给定的 A/B 顺序计能差。激发态使用 Tamm--Dancoff 近似，近简并态按事先规定的能量与强度容差匹配。频率题须在收敛几何排除平动和转动，以内部振动无虚频判定局部极小点，并按指定 C=O 键轴方向的相对位移归属伸缩模式。几何与方法、真实 ORCA 作业、收敛证据和题目交付物均满足，才判该次作答通过。

\textbf{四组分工。}人工组由我编写输入、运行 ORCA 并独立读数，第二位化学读者核对输入与原始输出后形成参考答案。固定脚本组禁用大语言模型，使用预先写好的参数化流程。脚本编写和题面转参的人工时间单列。专用 Auto-ORCA 组和加载冻结版 ORCA Skill 的通用智能体组各自规划、执行和解释。辅助脚本未覆盖的输入形式可经共同受控入口提交，所有启动均须留有回执。

人工与固定脚本每题各作答一次，两种智能体各在新会话中作答两次，交替先后顺序。共 144 次作答，其中 96 次为智能体会话。自动化组无法访问人工输入、答案键和其他组记录。评阅者隐藏组别后对照原始证据评分。已有 DFT 评测发现部分高分答案的运行记录不完整\cite{dftbench}，本研究因此要求每次作答提供可核对的作业证据。

\clearpage
\subhead{研究二\quad 已有输出能回答哪些后续问题？}
为 12 种中性单重态分子准备同一几何下的气相与 CPCM 水相 B3LYP/def2-SVP 输出。其中四种沿用研究一人工组经复核的八份输出，其余八种另行计算。对每种分子提出两问，在相互独立的会话中给出完全相同的文件。六对只交换能量相减方向。三对先问电子能差，再问标准水合自由能。其余三对隐去水相文件，分别询问已有的气相能量和无法由现有证据推出的溶剂效应。

两个 Auto-ORCA 条件都必须给出单位和原始输出位置，或指出还缺少什么。只有干预组被明确要求逐项核对物理量、几何、理论水平和溶剂。每问每组重复两次，共 96 场只读会话。要求重算仅记为请求。能差容差为 $0.1$~\kJmol{}，总能量容差为 $10^{-5}$~hartree。两问均与经复核的答案一致才算该对通过。每个条件每轮有 18 个可给数值的问题和 6 个证据不足的问题，分别统计漏答、可避免的重算请求及无依据报数。

\subhead{研究三\quad 在四次 DFT 查询内选择构象}
另取 10 种仅含 C、H、O 的中性柔性分子，每种至少三个可旋转键、最多八个重原子。按固定随机种子生成最多 200 个候选，经 MMFF94 弛豫后采用对称性校正的重原子均方根偏差去重（阈值 $0.35$~\AA），按生成顺序保留前八个。用 MMFF94 与 GFN2-xTB 在这八套相同坐标上排序。另用气相 $r^2$SCAN-3c 计算全部 80 个单点能，复核几何与收敛后对选择者隐藏参考表。一次查表消耗一次 DFT 查询额度。

两种固定规则分别选 MMFF94 排名前四的候选、或按同一排名优先覆盖不同扭转角家族，再依序补足。两种 Auto-ORCA 条件都看到低成本排名和家族标记，从 MMFF94 最低者出发，共选四轮。每轮都返回候选编号与剩余额度，只有反馈组还看到所选候选的 DFT 能量。每种分子每条固定规则运行一次，每个智能体条件独立重复三次，共 60 场模型会话。AutoDFT 在 VASP 任务中比较过固定与可调整计划，这也是本研究同时记录成功率和计算成本的依据\cite{autodft}。

若选中几何的最低能量距八个候选的参考最低值不超过 $0.01$~\kJmol{}，记为命中。另记录两者的能量差。未选满四次计为未命中，并报告实际查询数及已选到的最低能量。零次查询则将能量差记为缺失。

\subhead{如何比较结果}
研究一按题配对，先在每题内汇总智能体的两次作答，再对每类四题等权汇总。人工结果提供参考值及计时基线。三个自动化组分别报告全题通过率、失败原因、重试次数、人工分钟、总耗时、CPU 核时及模型用量。人工组也记录操作和计算时间。逐题保留实际输入、原始输出和读数出处。研究二和三同样先在分子内汇总，再报告配对差值及逐一剔除分子后的变化范围。主评测复用到研究二的四种分子不再计作新的独立样本。

\clearpage
\section{预期结果与意义}
四组比较将给出按题型划分的正确性和成本表。若固定脚本与两种智能体在单点能题上同样准确，选择流程时可直接比较准备时间与人工干预次数。对于频率和激发态题，错误还可以定位到选错结构、读错谱表或作出缺乏依据的归属。我预计，证据检查对物理量发生变化的问题帮助最大。两套低成本构象排名分歧较大时，逐次返回 DFT 能量的价值较大。结果将限定于选定分子、软件配置及一名人工操作者，构象结论仅涉及冻结几何集合内的电子能排序。

\section{时间安排}
\begin{tabularx}{\linewidth}{|P{0.27\linewidth}|c|c|Y|}
\hline
\rowcolor[HTML]{F2F2F2}\textbf{任务} & \textbf{开始周} & \textbf{结束周} & \textbf{产出} \\ \hline
选题与判分设计 & 1 & 2 & 24 道题的坐标、交付物、分层与参考容差 \\ \hline
预试并冻结 & 3 & 4 & 每类预试、人工参考流程、软件版本与预算 \\ \hline
四组正式评测 & 5 & 7 & 每周八题的人工参考与 40 份自动化作答 \\ \hline
输出复用与构象选择 & 8 & 9 & 12 对问题、10 套构象能量表和选择记录 \\ \hline
复核与配对分析 & 10 & 11 & 分题型结果、失败分类与成本比较 \\ \hline
撰写及汇报 & 12 & 12 & 研究论文、图表及展示材料 \\ \hline
\end{tabularx}

\section{资源、可行性与风险}
拟使用 ORCA~6.1.1、Python、RDKit、xTB、两套智能体环境及一位独立复核者。四组统一求解器版本、数值设置与时限，每次只执行一项单核作业，进程树内存上限 1~GiB，每进程 MaxCore 为 192~MB。若两种智能体支持同一底层模型，将固定相同版本和推理设置。否则记录差异，只在所测系统配置层面解释组间结果。

\textbf{预算。}主评测每遍的两类双单点题需 $8+8$ 次 ORCA 启动，另四类各需四次。按优化与频率合并提交，每遍 32 次，六遍共 192 次。研究二需 24 份输出，其中八份可复用。研究三需 80 次单点。正式首轮合计 $192+(24-8)+80=288$ 次。预留 48 次供预试、拆分作业及重试，总上限拟为 336 次。每次限 15 分钟，求解器累计运行时间上界为 84 小时。正式模型会话共 $96+96+60=252$ 场，每场最多八次模型调用。人工准备、复核、模型等待与计费另记。

第 4 周前以开发样例核对各类作业的内存、时间、模型额度和复核工时，再冻结题面、输入、脚本、Skill 与答案规则。与开发样例重复或无法生成足够候选的分子，只能在揭盲前按同类预设顺序替换。正式阶段的超时和失败保留在分母中。预算耗尽或服务中断时报告各组已覆盖的题数。独立复核者审阅的是计算证据。仅记录我本人的操作工时。答案键和凭据与智能体工作目录隔离，按许可共享题面、判分程序和可发布的科学数据。

\clearpage
\renewcommand{\bibsection}{\section{参考文献}}
\setlength{\bibsep}{0.65em}
\begin{thebibliography}{9}
\raggedright
\bibitem{jacobs}
P. F. Jacobs and R. Pollice. Developing large language models for quantum
chemistry simulation input generation. \emph{Digital Discovery} 4 (2025),
762--775. \doi{10.1039/D4DD00366G}.

\bibitem{qplanner}
R. Zhang, W. Gao and Y. Kurashige. Q-Planner: A harness system for quantum
chemistry agents. \emph{ACS Omega} (2026). \doi{10.1021/acsomega.6c05744}.

\bibitem{quntur}
J. B. P\'erez-S\'anchez, Y. Zou, J. A. Campos-Gonzalez-Angulo, et al.
El Agente Quntur: A research collaborator agent for quantum chemistry.
\emph{Chem} (2026), 103259. \doi{10.1016/j.chempr.2026.103259}.

\bibitem{censo}
S. Grimme, F. Bohle, A. Hansen, P. Pracht, S. Spicher and M. Stahn.
Efficient quantum chemical calculation of structure ensembles and free
energies for nonrigid molecules. \emph{Journal of Physical Chemistry A}
125 (2021), 4039--4054. \doi{10.1021/acs.jpca.1c00971}.

\bibitem{dftbench}
S. Liu, J. Yang, W. Wan, J. Han, J. Li and T. Deng. DFT-Agent Bench:
Scaffold effects and execution provenance in autonomous DFT.
\emph{arXiv} (2026), \href{https://arxiv.org/abs/2611.10629v1}{2611.10629v1}
（预印本）。

\bibitem{autodft}
P. Yang, Z. Zhang, Y. Li, et al. AutoDFT: A closed-loop multi-agent framework
for autonomous DFT calculations. \emph{arXiv} (2026),
\href{https://arxiv.org/abs/2605.26179v2}{2605.26179v2} （预印本）。
\end{thebibliography}
\end{document}
